%% ========================================================================
%% CAPÍTULO 4 - CONCLUSÕES
%% ========================================================================

\capitulo{conclusao}{Conclusões}

Este trabalho apresentou uma análise abrangente da mobilidade urbana na Região Metropolitana de Brasília, com foco na sustentabilidade dos sistemas de transporte. A pesquisa combinou análise quantitativa de dados secundários, pesquisa de campo com usuários e análise qualitativa de políticas públicas, permitindo uma compreensão multidimensional dos desafios e oportunidades para o desenvolvimento de um sistema de mobilidade mais sustentável.

\section{Síntese dos Principais Resultados}

Os resultados obtidos confirmaram as hipóteses iniciais da pesquisa e revelaram aspectos importantes sobre os padrões de mobilidade na região:

\subsection{Caracterização dos Padrões de Mobilidade}

A análise dos padrões de deslocamento evidenciou a predominância do transporte individual motorizado, que representa 65,2\% dos deslocamentos na região. Esta proporção é significativamente superior à observada em cidades com sistemas de transporte mais desenvolvidos, confirmando a dependência excessiva do automóvel particular.

O transporte público atende apenas 24,8\% da demanda de deslocamentos, percentual considerado baixo para uma região metropolitana de mais de 3 milhões de habitantes. Os modos não motorizados (caminhada e bicicleta) representam apenas 9,4\% dos deslocamentos, indicando potencial significativo de crescimento com investimentos adequados em infraestrutura.

\subsection{Qualidade dos Serviços de Transporte}

A avaliação da qualidade do transporte público revelou deficiências significativas em todas as dimensões analisadas. A média geral de 2,2 pontos (em escala de 5) classifica o serviço como "ruim", com problemas especialmente graves relacionados à integração entre modais (1,9 pontos) e conforto dos veículos (2,0 pontos).

Os tempos de deslocamento no transporte público (78 minutos em média) são mais que o dobro dos observados no transporte individual (35 minutos), constituindo uma barreira importante para a migração modal. Esta diferença está relacionada à baixa frequência dos serviços, necessidade de múltiplas transferências e infraestrutura inadequada.

\subsection{Impactos Ambientais e Socioeconômicos}

A análise dos impactos ambientais mostrou que o modelo atual de mobilidade gera emissões elevadas de gases poluentes. O automóvel particular emite 120 g CO₂/pass-km, valor 2,7 vezes superior ao ônibus e 4,8 vezes superior ao metrô. A predominância deste modal resulta em impactos ambientais significativos que poderiam ser reduzidos com a migração para modos mais sustentáveis.

Os custos socioeconômicos totais do transporte individual (R\$ 0,95/pass-km) são mais que o dobro dos custos do transporte público (R\$ 0,45/pass-km), evidenciando a ineficiência econômica do modelo atual. Estes custos incluem não apenas os gastos diretos dos usuários, mas também os custos sociais relacionados a congestionamentos, acidentes e poluição.

\section{Contribuições da Pesquisa}

Esta pesquisa oferece contribuições importantes em três dimensões:

\subsection{Contribuições Metodológicas}

O trabalho desenvolveu uma metodologia integrada para avaliação da mobilidade urbana que combina análise quantitativa de dados secundários, pesquisa de campo e análise qualitativa. Esta abordagem mostrou-se eficaz para compreender a complexidade dos fenômenos relacionados à mobilidade urbana e pode ser replicada em outras regiões metropolitanas.

O conjunto de indicadores proposto para monitoramento das políticas de mobilidade urbana constitui uma ferramenta prática para gestores públicos, permitindo acompanhar o progresso das intervenções e realizar ajustes necessários.

\subsection{Contribuições Empíricas}

A pesquisa fornece um diagnóstico detalhado e atualizado da situação da mobilidade urbana na Região Metropolitana de Brasília, preenchendo uma lacuna importante na literatura sobre o tema. Os dados coletados e analisados constituem uma base de referência para estudos futuros e para o desenvolvimento de políticas públicas.

A identificação das principais barreiras para adoção de modos de transporte mais sustentáveis oferece subsídios importantes para o direcionamento de investimentos e políticas públicas. Os resultados mostram que as deficiências do transporte público são o principal obstáculo para a migração modal.

\subsection{Contribuições Práticas}

As estratégias integradas propostas para promoção da mobilidade sustentável constituem um plano de ação abrangente e factível para a região. As propostas estão fundamentadas em evidências empíricas e seguem as melhores práticas internacionais, aumentando as chances de sucesso na implementação.

O trabalho oferece recomendações específicas para diferentes atores (governo federal, estadual, municipal e setor privado), facilitando a coordenação de esforços e a implementação das estratégias propostas.

\section{Limitações do Estudo}

É importante reconhecer as limitações desta pesquisa:

\subsection{Limitações Temporais}

O período de análise (2020-2023) foi influenciado pela pandemia de COVID-19, que alterou significativamente os padrões de mobilidade urbana. Embora tenham sido feitos ajustes metodológicos para considerar estes efeitos, alguns resultados podem não refletir completamente as condições normais de mobilidade.

\subsection{Limitações Espaciais}

O foco na Região Metropolitana de Brasília, embora justificado pelas características específicas da região, limita a generalização dos resultados para outras regiões metropolitanas brasileiras. Estudos comparativos seriam necessários para validar a aplicabilidade das estratégias propostas em outros contextos.

\subsection{Limitações Metodológicas}

O uso de questionários para coleta de dados primários pode introduzir viés de resposta, especialmente relacionado à desejabilidade social. Métodos complementares, como observação direta e análise de dados de GPS, poderiam enriquecer a análise.

\section{Recomendações para Trabalhos Futuros}

Com base nos resultados obtidos e nas limitações identificadas, recomenda-se:

\subsection{Estudos Longitudinais}

Realizar estudos longitudinais para acompanhar a evolução dos padrões de mobilidade ao longo do tempo, especialmente após a implementação das estratégias propostas. Estes estudos permitiriam avaliar a efetividade das intervenções e identificar tendências de longo prazo.

\subsection{Estudos Comparativos}

Desenvolver estudos comparativos entre diferentes regiões metropolitanas brasileiras, identificando fatores que influenciam o sucesso de políticas de mobilidade sustentável. Esta abordagem permitiria identificar boas práticas e adaptar estratégias para diferentes contextos.

\subsection{Análise de Viabilidade Econômica}

Realizar análises detalhadas de viabilidade econômica das estratégias propostas, incluindo análise custo-benefício e identificação de fontes de financiamento. Estes estudos são essenciais para subsidiar decisões de investimento e priorização de projetos.

\subsection{Impactos da Tecnologia}

Investigar os impactos de novas tecnologias (veículos elétricos, sistemas inteligentes de transporte, aplicativos de mobilidade) nos padrões de mobilidade urbana. A rápida evolução tecnológica pode alterar significativamente os cenários futuros de mobilidade.

\subsection{Participação Social}

Desenvolver estudos sobre participação social em políticas de mobilidade urbana, investigando mecanismos efetivos para envolver a população nas decisões sobre transporte. A aceitação social é fundamental para o sucesso de políticas de mobilidade sustentável.

\section{Implicações para Políticas Públicas}

Os resultados desta pesquisa têm implicações importantes para o desenvolvimento de políticas públicas de mobilidade urbana:

\subsection{Priorização de Investimentos}

Os resultados indicam que os investimentos devem priorizar a melhoria da qualidade do transporte público, especialmente em aspectos como frequência, pontualidade e integração entre modais. Estes investimentos têm potencial de gerar impactos significativos na migração modal.

\subsection{Coordenação Institucional}

A complexidade dos desafios de mobilidade urbana requer coordenação entre diferentes níveis de governo e setores. A criação de uma autoridade metropolitana de transporte poderia facilitar esta coordenação e melhorar a eficiência das políticas implementadas.

\subsection{Financiamento Sustentável}

O desenvolvimento de um sistema de mobilidade sustentável requer fontes de financiamento diversificadas e sustentáveis. A combinação de recursos públicos, parcerias público-privadas e mecanismos inovadores de financiamento é essencial para viabilizar os investimentos necessários.

\section{Considerações Finais}

A transição para um sistema de mobilidade urbana mais sustentável na Região Metropolitana de Brasília é um desafio complexo que requer ação coordenada de múltiplos atores. Os resultados desta pesquisa mostram que, apesar das dificuldades atuais, existem oportunidades significativas para melhorar a sustentabilidade do sistema de transporte.

As estratégias propostas, se implementadas de forma integrada e coordenada, têm potencial para transformar significativamente os padrões de mobilidade na região. No entanto, o sucesso dessas estratégias depende de vontade política, recursos financeiros adequados e participação ativa da sociedade.

A mobilidade urbana sustentável não é apenas uma questão técnica, mas também social, econômica e ambiental. Sua promoção requer uma visão sistêmica que considere as interações entre transporte, uso do solo, desenvolvimento econômico e qualidade de vida urbana.

Este trabalho contribui para o avanço do conhecimento sobre mobilidade urbana sustentável e oferece subsídios práticos para gestores públicos e pesquisadores. Espera-se que os resultados apresentados possam apoiar o desenvolvimento de políticas mais efetivas e a construção de cidades mais sustentáveis e equitativas.

A jornada em direção à mobilidade urbana sustentável é longa e desafiadora, mas os benefícios potenciais - em termos de qualidade de vida, sustentabilidade ambiental e desenvolvimento econômico - justificam plenamente os esforços necessários. O momento é propício para ação, e esta pesquisa oferece um roteiro para o caminho a seguir.